\documentclass[11pt, fleqn]{amsart}
\usepackage{amsmath, amssymb, amsthm} 
\usepackage{hyperref}
\usepackage{mathtools}
\usepackage{graphicx}
\usepackage{adjustbox}
\usepackage{caption}
\usepackage{subcaption}
\usepackage{booktabs}
\usepackage{float}
\usepackage{longtable}
\usepackage{tikz-cd}
% Code listing packages (for the Key to Tables section)
\usepackage{listings}
\usepackage{xcolor}
\lstset{
  basicstyle=\ttfamily\small,
  breaklines=true,
  frame=single,
  language=Python,
  keywordstyle=\color{blue},
  commentstyle=\color{gray},
  stringstyle=\color{red}}
% Theorem environment setup
\newtheorem{theorem}{Theorem}[section]
\newtheorem{lemma}[theorem]{Lemma}
\newtheorem{corollary}[theorem]{Corollary}
\newtheorem{conjecture}{Conjecture}
\newtheorem{proposition}{Proposition}
\newtheorem*{proposition*}{Proposition}
\newtheorem{observation}{Observation}
\theoremstyle{remark}
\newtheorem{remark}{Remark}
% Title and authorship
\title
{Details of the undeformed $h=2$ argument: a non-oscillating companion}
\author{Barry Brent}
\email{axcjh@bu.edu}
\subjclass[2020]{05A17, 05C62, 11F11, 11P83}
\keywords{partitions, graphs, cusp forms}
\begin{document}
\maketitle
\vspace{-0.5em}
\begin{center}
    %rm
    {14 June 2026}
\end{center}
\vspace{1em}
Say $h_0 = 1$ and $h_n = 2$ for $n$ positive. This is the companion of the deformed
calculation: here we make \emph{no} deformation, introduce no parameter $c$, and study
the coefficients $h_n$ themselves. From Lemma 2.1, with
$H(t) = \sum_{n\ge 0} h_n t^n$ and $J(t) = \sum_{r\ge 1} j_r t^r$,
the following statements are equivalent:
\newline
\newline
(i) $
  t H'(t) = H(t)\,J(t),
$ \newline
(ii) $
  n h_n = \sum_{r=1}^{n} j_r\, h_{n-r}\ \ (n \ge 1).
$ \newline \newline
From (a), we also have \newline
(iii)
$
  J(t) = t\,\frac{d}{dt}\log H(t).
$
\newline
\newline
With $h_0 = 1$ and $h_r = 2$ for $r \ge 1$, summing the geometric series gives the
closed form
$$
  H(t) = 1 + 2\sum_{r\ge 1} t^r = 1 + \frac{2t}{1-t} = \frac{1+t}{1-t}.
$$
This is the object of study below. For comparison with the deformed family we record
that
$$
  \frac{d}{dt}\log H(t)
    = \frac{d}{dt}\bigl(\log(1+t) - \log(1-t)\bigr)
    = \frac{1}{1+t} + \frac{1}{1-t}
    = \frac{2}{1-t^2},
$$
so $J(t) = t\cdot\dfrac{2}{1-t^2} = 2t + 2t^3 + 2t^5 + \cdots$, i.e.\ $j_r = 2$ for odd
$r$ and $j_r = 0$ for even $r$. The deformed analysis prepends a parameter $c$ to this
$j$-sequence and arrives at $H^{(c)}(t)=e^{ct}/(1-t^2)$; the present note keeps $J$ as
it is and analyses $H(t)=(1+t)/(1-t)$ directly. We will see that the two differ in
exactly one structural respect --- the singularity at $t=-1$ --- and that this single
difference is the whole of the oscillation phenomenon.

\medskip

We give a self-contained analytic derivation of the large-$n$ behavior of the
coefficients $h_n$ of the undeformed generating function $H(t)=(1+t)/(1-t)$. The only
input is the closed form on the unit disc; no individual coefficient is computed and no
numerical evidence is invoked. We obtain an exact formula for $h_n$, show that the
sequence is \emph{eventually exactly constant} rather than merely convergent, and obtain
the single limit $2$ with no parity splitting. A closing section compares this with the
deformed family $H^{(c)}(t)=e^{ct}/(1-t^2)$, whose coefficients split into $\cosh c$ and
$\sinh c$, and isolates the simple pole at $t=-1$ as the unique source of that
oscillation.

\section{Setup and notation}
We take as our single starting point the identity, valid as an equality of analytic
functions on the open unit disc $\{\,t\in\mathbb{C}:|t|<1\,\}$,
\begin{equation}
H(t)\;=\;\frac{1+t}{1-t}\;=\;\sum_{n=0}^{\infty}h_n\,t^n.
\tag{1}
\end{equation}
This is what equation~(D) produces for the constant sequence $h_0=1,\ h_n=2\ (n\ge1)$
\emph{before any deformation}: summing the geometric series gives the rational function
on the left, and the two sides, being analytic on the unit disc, agree throughout it.
The numbers $h_n$ are exactly the Taylor coefficients in $(1)$; their large-$n$ behavior
is the object of study. Nothing below uses the value of any particular $h_n$.

Throughout, ``boundary'' refers to the circle $|t|=1$, the boundary of the disc of
convergence in $(1)$. The only point of $\overline{\{|t|<1\}}$ at which $H$ fails to be
analytic is $t=+1$; the point $t=-1$, where the denominator $1-t$ does not vanish, is
not a singularity. This asymmetry is the entire content of what follows.

\section{An exact formula for the coefficients}

\begin{lemma}\label{lem:exact}
For every $n\ge 0$,
\[
h_n\;=\;2-\delta_{n,0}
\qquad\text{equivalently}\qquad
h_0=1,\quad h_n=2\ \ (n\ge1),
\]
where $\delta_{n,0}$ is the Kronecker delta.
\end{lemma}

\begin{proof}
We decompose $H$ over the two boundary points $t=\pm1$, exactly as the deformed
calculation decomposes $1/(1-t^2)$. Since $H$ is a ratio of polynomials of equal degree,
the decomposition carries a polynomial part:
\[
\frac{1+t}{1-t}
\;=\;\underbrace{\frac{2}{1-t}}_{\text{pole at }t=+1}
\;+\;\underbrace{\frac{0}{1+t}}_{\text{pole at }t=-1}
\;-\;1 ,
\]
an equality of rational functions on $\mathbb{C}\setminus\{1\}$; one checks it by
clearing denominators, $\,2-(1-t)=1+t$. The coefficient attached to the second boundary
point $t=-1$ is identically zero. On the disc $|t|<1$ the surviving pole expands as the
absolutely convergent geometric series $\frac{1}{1-t}=\sum_{i\ge0}t^{i}$, so that
\[
[t^n]\,\frac{2}{1-t}=2\quad(n\ge0),
\qquad
[t^n]\,(-1)=-\delta_{n,0},
\qquad
[t^n]\,\frac{0}{1+t}=0 .
\]
Since the Taylor coefficients of the analytic function $H$ are unique, comparison with
$(1)$ gives $h_n=2-\delta_{n,0}$.
\end{proof}

\noindent
It is worth isolating what the proof used. The expansion $\frac{1}{1+t}=\sum_{i\ge0}(-1)^i t^i$
is the only place a factor $(-1)^n$ could enter a coefficient of $H$; it is attached to
the pole at $t=-1$. Here that pole is absent (its coefficient is $0$), so no $(-1)^n$
ever appears. In the deformed companion the same expansion is present with coefficient
$\tfrac12 e^{-c}\neq0$, and it is precisely that term which alternates.

\section{A termination estimate (the analogue of the exponential tail)}

In the deformed problem the numerator $e^{ct}$ is an infinite series, and the error in
the limiting expression is governed by the tail of that series. Here the numerator
$1+t$ is a polynomial, the convolution terminates, and the corresponding ``tail'' is not
merely small but identically zero past the degree of the numerator. We state this in the
form we need.

\begin{lemma}\label{lem:terminate}
Let $P(t)=\sum_{k=0}^{d}p_k t^k$ be a polynomial and $H(t)=P(t)/(1-t)=\sum_{n\ge0}h_n t^n$
on $|t|<1$. Then for every $n\ge d$,
\[
h_n\;=\;P(1)\;=\;\sum_{k=0}^{d}p_k ,
\]
a constant independent of $n$. In particular the deviation of $h_n$ from this constant
is exactly $0$ for all $n\ge d$.
\end{lemma}

\begin{proof}
On $|t|<1$ write $H(t)=P(t)\sum_{m\ge0}t^m$, a product of an entire polynomial and an
absolutely convergent geometric series; the Cauchy product converges absolutely and its
$t^n$ coefficient is the finite convolution
\[
h_n=\sum_{k=0}^{\min(n,d)}p_k .
\]
For $n\ge d$ the upper limit is $d$, so $h_n=\sum_{k=0}^{d}p_k=P(1)$, independent of $n$.
\end{proof}

\begin{corollary}\label{cor:terminate2}
For $H(t)=(1+t)/(1-t)$ we have $P(t)=1+t$, so $d=1$ and $P(1)=2$; hence $h_n=2$ for all
$n\ge1$. Writing $\varepsilon_n:=h_n-2$ for the deviation from the eventual value, we
have
\[
\varepsilon_n=-\delta_{n,0},\qquad\text{so}\qquad \varepsilon_n=0\ \text{for all } n\ge1 .
\]
\end{corollary}

\noindent
This is the structural counterpart of the deformed tail bound
$|\varepsilon_n(c)|\le |c|^{\,n+1}/(n+1)!\,e^{|c|}$. There the error decays
super-exponentially but is positive for every $n$; here it is not merely small, it is
zero. The reason is the same in both problems --- the size of the error is controlled by
how much of the numerator's expansion has yet to enter the convolution --- but a
polynomial numerator exhausts itself after finitely many steps, whereas an exponential
numerator never does.

\section{The asymptotics}

\begin{theorem}\label{thm:main}
For every $n\ge0$,
\[
h_n\;=\;2\;+\;\varepsilon_n,\qquad \varepsilon_n=-\delta_{n,0},
\]
so $\varepsilon_n=0$ for all $n\ge1$. Consequently the full sequence converges,
\[
\lim_{n\to\infty}h_n=2 ,
\]
and the even and odd subsequences share this single limit,
\[
\lim_{\substack{n\to\infty\\ n\ \mathrm{even}}}h_n
=\lim_{\substack{n\to\infty\\ n\ \mathrm{odd}}}h_n
=2 .
\]
There is no parity splitting: the limit does not depend on the parity of $n$.
\end{theorem}

\begin{proof}
The exact formula is Lemma~\ref{lem:exact}, equivalently
Corollary~\ref{cor:terminate2} with $\varepsilon_n=h_n-2=-\delta_{n,0}$. For $n\ge1$ we
have $\varepsilon_n=0$, so $h_n=2$ identically; a sequence equal to $2$ from $n=1$
onward converges to $2$, and every subsequence --- in particular the even-$n$ and odd-$n$
subsequences --- has the same limit $2$. Since no factor $(-1)^n$ occurs (the pole at
$t=-1$ being absent, by Lemma~\ref{lem:exact}), there is nothing to distinguish the two
parities.
\end{proof}

\begin{corollary}[contrast with the deformed family]\label{cor:contrast}
In the deformed companion, for $H^{(c)}(t)=e^{ct}/(1-t^2)$ one has
$h^{(c)}_n=\tfrac12\bigl(e^{c}+(-1)^n e^{-c}\bigr)+\varepsilon_n(c)$, with even-$n$ limit
$\cosh c$ and odd-$n$ limit $\sinh c$, two values that coincide only when $\sinh c=0$,
i.e.\ $c=0$. By contrast the undeformed sequence has a single limit $2$ for both
parities, for every realisation of the construction. The deformed even and odd limits do
not specialise to the undeformed value: as $c\to0$ they tend to $1$ and $0$, the
coefficients of $1/(1-t^2)$, not to $2$. The undeformed sequence $(1+t)/(1-t)$ is
therefore not a member of the one-parameter deformed family $e^{ct}/(1-t^2)$, and the
comparison below is between two genuinely different generating functions, not two values
of one parameter.
\end{corollary}

\section{Comparison with the deformed family: locating the oscillation}

\begin{remark}[which pole supplies what]\label{rem:poles}
Place the two generating functions side by side on the closed unit disc and list their
singularities on the boundary circle $|t|=1$:
\[
H^{(c)}(t)=\frac{e^{ct}}{(1-t)(1+t)}\quad\text{(poles at }t=+1\text{ and }t=-1),
\qquad
H(t)=\frac{1+t}{1-t}\quad\text{(pole at }t=+1\text{ only).}
\]
In each case the coefficient asymptotics are governed by these boundary poles. A simple
pole at $t=+1$ contributes a constant (parity-independent) mode to $h_n$; a simple pole
at $t=-1$ contributes the mode $(-1)^n$, because $\frac{1}{1+t}=\sum_n(-1)^n t^n$.

For the deformed function both poles are present and neither is cancelled, since the
numerator $e^{ct}$ is entire and zero-free. The pole at $t=+1$ supplies the constant
$\tfrac12 e^{c}$ and the pole at $t=-1$ supplies the alternating $\tfrac12(-1)^n e^{-c}$;
their sum is $\cosh c$ on even $n$ and $\sinh c$ on odd $n$. The splitting of the limit
into two distinct values --- the oscillation --- is produced \emph{solely} by the
\emph{location} $t=-1$ of the second pole, and the parameter $c$ only fixes the two
residue magnitudes $e^{\pm c}/2$.

For the undeformed function the numerator $1+t$ vanishes at $t=-1$ \emph{exactly where
the second pole would sit}, cancelling it. Only the pole at $t=+1$ survives, only the
constant mode is present, and the alternating mode has coefficient $0$
(Lemma~\ref{lem:exact}). Hence a single limit and no oscillation. The contrast between
oscillation and its absence is thus, quite literally, the presence or absence of a
singularity at the single point $t=-1$.
\end{remark}

\begin{remark}[a one-line criterion]\label{rem:criterion}
The mechanism is general. If $G$ is meromorphic on a neighborhood of $\overline{\{|t|<1\}}$
with its only boundary singularities simple poles at points $e^{i\theta_1},\dots,e^{i\theta_k}$,
then $[t^n]G$ is asymptotically a sum of terms $\rho_j\,e^{-i n\theta_j}$, one per pole.
The sequence converges to a single limit (no oscillation) precisely when the only
boundary pole sits at $\theta=0$, i.e.\ at $t=+1$; any boundary pole with $\theta\neq0$
contributes a genuinely oscillating mode of angular frequency $\theta$. The value
$\theta=\pi$, i.e.\ a pole at $t=-1$, gives the period-$2$ mode $(-1)^n$ and hence the
even/odd dichotomy seen in the deformed coefficients. In one sentence:
\emph{oscillation in the coefficient sequence is the signature of a dominant boundary
singularity off the positive real axis.} The deformation installs such a singularity at
$t=-1$ (via the factor $1-t^2$, which the zero-free numerator $e^{ct}$ cannot repair);
the undeformed sequence removes it (via the numerator zero of $1+t$ at $t=-1$).
\end{remark}

\begin{remark}[relation to the spectral plots]\label{rem:spectral}
The oscillation studied empirically in the main article lives one level further on: it
is in the minimum-modulus eigenvalues $\mu_n$ of the matrices $J^{(c)}_{h,n}$, not
directly in the coefficients $h_n$. The present note does not prove anything about
$\mu_n$. What it does supply is the analytic shadow of that phenomenon at the level of
the defining sequence: a period-$2$ ($(-1)^n$) structure is present in the deformed
coefficients and absent in the undeformed ones, for exactly the reason given in
Remark~\ref{rem:poles}. This is consistent with the observation recorded for the
$2$-sequence (the deformed $c=1$ plot, with its even subsequence approaching a constant
and its odd subsequence sitting at a separate value), and with the flatness of the
undeformed picture. We flag the link as heuristic: whether the boundary-pole dichotomy
propagates through the determinant and characteristic-polynomial machinery to control the
$\mu_n$ is a separate question, and is not settled here.
\end{remark}

\begin{remark}[independence from numerics]\label{rem:numerics}
Every assertion follows from $(1)$ together with the elementary
Lemmas~\ref{lem:exact} and~\ref{lem:terminate}. In particular the conclusion ``no
oscillation'' is not an extrapolation from computed coefficients: for each $n\ge1$ the
coefficient $h_n$ is shown to equal $2$ exactly, and the deviation $\varepsilon_n$ is
shown to vanish identically, with no appeal to a table of values.
\end{remark}

\end{document}
